%% ============================================================
%%  CHAPTER 12 — REAL-TIME INFERENCE
%% ============================================================
\chapter{Real-Time Inference}
\label{ch:inference}

This chapter describes how the offline-trained \acrshort{dtw}
hybrid model (Chapter~9) is deployed as a continuously running,
real-time classifier on the Raspberry Pi~5 (Chapter~4), turning a
live \acrshort{eeg} stream into discrete prosthetic hand commands
sent over \acrshort{uart} to the ESP32 (Chapter~11). Every
component in this chapter reuses the exact model, features, and
preprocessing fitted offline in Chapters~7--9; nothing is retrained
or re-derived for deployment. Section~\ref{sec:inf_overview}
presents the end-to-end real-time pipeline. Section~\ref{sec:inf_window}
describes the sliding window acquisition strategy and how it
differs from offline epoching. Section~\ref{sec:inf_gate}
describes the real-time activity gate and how it differs from the
recording-relative gate used offline (Chapter~7).
Section~\ref{sec:inf_classify} describes per-window classification.
Section~\ref{sec:inf_confidence} describes confidence thresholding.
Section~\ref{sec:inf_fsm} describes the Finite State Machine that
confirms a command before it is acted on.
Section~\ref{sec:inf_latency} presents the latency budget.
Section~\ref{sec:inf_uart} describes the \acrshort{uart} output
protocol (Contract~A, Chapter~11) and corrects an outdated
description found in earlier project notes.
Section~\ref{sec:inf_compare} summarises every offline/real-time
pipeline difference in one place. Section~\ref{sec:inf_summary}
summarises the chapter.

%% ----------------------------------------------------------
\section{Real-Time Pipeline Overview}
\label{sec:inf_overview}
%% ----------------------------------------------------------

Figure~\ref{fig:realtime_pipeline} shows the complete real-time
pipeline, from the continuous \acrshort{eeg} stream to the
\acrshort{uart} line sent to the ESP32. The pipeline runs
continuously on the Raspberry Pi~5, re-evaluating a new window
every 0.2\,s, independently of whether the previous window produced
a command.

\begin{figure}[htbp]
\centering
\begin{tikzpicture}[
    node distance=0.5cm,
    box/.style={rectangle, rounded corners=3pt, draw=black,
                fill=blue!10, minimum width=6.2cm,
                minimum height=0.75cm, align=center, font=\small},
    decision/.style={diamond, draw=black, fill=yellow!20,
                     minimum width=3.4cm, minimum height=0.9cm,
                     align=center, font=\small},
    arrow/.style={-{Stealth[length=5pt]}, thick},
    note/.style={font=\scriptsize\itshape, text=gray, align=left}
]

\node[box, fill=green!15] (stream) {Continuous EEG Stream (19\,ch, 200\,Hz)};
\node[box, below=of stream] (window) {Sliding Window (2.0\,s, step 0.2\,s)};
\node[decision, below=of window] (gate) {Activity $\geq$ 20\,\si{\micro\volt}?};
\node[box, right=2cm of gate, fill=red!15] (skip) {No command\\(hand holds position)};
\node[box, below=of gate] (route) {Stage 1: blink vs.\ muscle};
\node[box, below=of route] (predict) {Stage 2a (DTW) or Stage 2b (Riemannian)};
\node[decision, below=of predict] (conf) {Confidence $\geq$ 0.60?};
\node[box, below=of conf] (fsm) {FSM: 3 consecutive identical predictions};
\node[box, below=of fsm, fill=orange!20] (uart)
    {UART line to ESP32 (Contract A, Chapter~11)};

\draw[arrow] (stream)  -- (window);
\draw[arrow] (window)  -- (gate);
\draw[arrow] (gate)    -- node[right,font=\scriptsize]{Yes} (route);
\draw[arrow] (gate)    -- node[above,font=\scriptsize]{No}  (skip);
\draw[arrow] (route)   -- (predict);
\draw[arrow] (predict) -- (conf);
\draw[arrow] (conf)    -- node[right,font=\scriptsize]{Yes} (fsm);
\draw[arrow] (conf)    -| node[near start, above,
    font=\scriptsize]{No} (skip);
\draw[arrow] (fsm)     -- (uart);

\node[note, right=0.3cm of window] {inference.py: 400-sample buffer, 40-sample step};
\node[note, right=0.3cm of gate]   {Fixed threshold (Section~\ref{sec:inf_gate})};
\node[note, right=0.3cm of predict] {Chapter~9 model, unchanged};
\node[note, right=0.3cm of fsm]     {Section~\ref{sec:inf_fsm}};

\end{tikzpicture}
\caption{The real-time inference pipeline. A new window is
evaluated every 0.2\,s; a window that fails the activity gate or
the confidence threshold, or does not yet complete an FSM run of 3
consecutive identical predictions, produces no UART output and the
hand simply holds its current position (Chapter~11's fail-safe
behaviour handles the case where no command arrives for an extended
period).}
\label{fig:realtime_pipeline}
\end{figure}

The design principle underlying every stage in this pipeline is the
same one established in Chapter~9: \textbf{nothing about the model
changes between offline evaluation and real-time deployment.} The
same 129-feature superset (Chapter~8), the same Stage~1 router, the
same \acrshort{dtw}-KNN blink classifier, and the same Riemannian
muscle classifier (all Chapter~9) are loaded from
\code{dtw\_hybrid\_model.pkl} and applied unmodified. Everything
described in this chapter is deployment \emph{machinery} around
that fixed model --- window buffering, gating, confidence
thresholding, and command confirmation --- not a re-implementation
of the model itself.

%% ----------------------------------------------------------
\section{Sliding Window Acquisition}
\label{sec:inf_window}
%% ----------------------------------------------------------

The real-time pipeline maintains a rolling buffer of the most
recent 2.0\,s of \acrshort{eeg} samples (400 samples at
\SI{200}{\hertz}, exactly matching the offline epoch length,
Chapter~7) and re-evaluates a classification every \textbf{0.2\,s}
(a 40-sample step), i.e.\ five predictions per second.

\subsection{Why the Real-Time Step Differs from the Offline Step}

Offline, epochs are extracted with a 0.5\,s step (75\% overlap,
Chapter~7) --- a choice driven by wanting many diverse, overlapping
training examples per recording, with the resulting overlap-leakage
risk handled at the evaluation-protocol level (Chapter~9) rather
than by reducing the step. In real-time deployment, the same
consideration does not apply --- there is no train/test split to
leak across, since inference makes a fresh prediction on live data
it was never trained on. The real-time step is instead chosen
purely to balance responsiveness against redundant computation: a
0.2\,s step gives five prediction opportunities per second, which
comfortably supports the \acrshort{fsm}'s 3-consecutive-prediction
confirmation requirement (Section~\ref{sec:inf_fsm}) reaching a
decision within roughly 0.6--1.0\,s of a gesture's onset, without
re-computing a full window more often than the underlying 14\,ms
per-window latency (Section~\ref{sec:inf_latency}) would comfortably
support at a much finer step.

%% ----------------------------------------------------------
\section{Real-Time Activity Gate}
\label{sec:inf_gate}
%% ----------------------------------------------------------

\subsection{A Different Gate From Training}

Chapter~7 established that the \emph{offline} activity gate used to
build the training dataset is \textbf{recording-relative}:
\code{gate\_mask(frac=0.5)} thresholds each epoch against 50\% of
that specific recording's own 90th-percentile activity level. This
works offline because the entire recording is available in advance
to compute a percentile from. It does \textbf{not} work in real
time: a streaming pipeline has no access to ``the rest of the
recording'' to compute a percentile over, since the recording is
still ongoing and, in normal operation, has no fixed end at all.

The real-time pipeline therefore uses a second, independently
calibrated gate: a \textbf{fixed} \acrshort{rms} threshold of
20\,\si{\micro\volt} at Fp1/Fp2 and T3/T4 (the same channel pairing
used by the offline gate's activity measure, Equation~7.6),
calibrated empirically against the same rest/active amplitude
distributions documented in Chapter~7:

\begin{itemize}[leftmargin=*]
    \item \textbf{Rest windows:} typically
    7--13\,\si{\micro\volt} RMS (background \acrshort{eeg}).
    \item \textbf{Active gesture windows:} typically
    30--200\,\si{\micro\volt} RMS.
\end{itemize}

A fixed threshold of 20\,\si{\micro\volt} sits comfortably in the
gap between these two distributions. Unlike the offline gate, this
threshold cannot adapt to a specific recording's own scale, so it
depends on electrode impedance and amplifier gain being
reasonably consistent with the conditions under which it was
calibrated --- a limitation carried forward to Chapter~14 as a
future-work item (re-calibrating this threshold per session, or
replacing it with a short rolling-window percentile that
approximates the offline gate's adaptivity without requiring the
full recording in advance).

\subsection{Why Two Gates, Not One}

It is tempting to ask why the same recording-relative gate is not
simply used in both places. The reason is architectural, not a
missed simplification: the recording-relative gate's entire value
proposition is that it adapts to a specific recording's own
amplitude scale, which requires that recording's full activity
distribution to already be known. That precondition is true for a
completed training recording and false for a live stream. The two
gates therefore serve the same conceptual purpose (reject epochs
that are almost certainly rest, not gesture) through two different,
context-appropriate mechanisms, and should not be confused with one
another --- a point made explicitly in Chapters~7 and~10 and
restated here in full for the real-time context.

%% ----------------------------------------------------------
\section{Per-Window Classification}
\label{sec:inf_classify}
%% ----------------------------------------------------------

A window that passes the activity gate is classified using exactly
the production pipeline described in Chapter~9:

\begin{enumerate}[leftmargin=*]
    \item The 129-feature superset (Chapter~8) is extracted from the
    window.
    \item Stage~1 routes the window to either the blink branch or
    the muscle branch based on frontal and temporal features.
    \item If routed to blink, the window's z-scored frontal
    waveform is classified by the \acrshort{dtw}-KNN sub-classifier
    (Chapter~9, Section~9.6) into \code{single\_blink} or
    \code{double\_blink}.
    \item If routed to muscle, the window's raw covariance is
    projected into Riemannian tangent space and classified
    (Chapter~9, Section~9.5) into \code{clinch},
    \code{bruxism\_left}, or \code{bruxism\_right}.
\end{enumerate}

Each sub-classifier returns both a predicted class and a confidence
score (class probability for the Riemannian branch's \acrshort{svc};
neighbour-vote agreement fraction for the \acrshort{dtw}-KNN
branch), which feeds directly into the confidence gate
(Section~\ref{sec:inf_confidence}).

%% ----------------------------------------------------------
\section{Confidence Thresholding}
\label{sec:inf_confidence}
%% ----------------------------------------------------------

A prediction is discarded --- treated identically to a window that
failed the activity gate --- unless its confidence score meets or
exceeds \textbf{0.60}. This threshold exists to suppress low-margin
predictions near a decision boundary from ever reaching the
\acrshort{fsm} (Section~\ref{sec:inf_fsm}) at all: a window the
model is genuinely unsure about is more useful to the system as ``no
command'' than as a low-confidence guess that might spuriously
begin accumulating toward an incorrect \acrshort{fsm} confirmation.
Combined with the \acrshort{fsm}'s 3-consecutive-prediction
requirement, this makes the system doubly conservative about acting
on a single ambiguous window --- consistent with the fail-closed
design philosophy already established for Contract~A
(Chapter~11): a missed or delayed command is a minor usability
issue, while an incorrectly executed one, on a prosthetic hand near
the user's face and body, is a safety concern.

%% ----------------------------------------------------------
\section{Finite State Machine Confirmation}
\label{sec:inf_fsm}
%% ----------------------------------------------------------

\subsection{Motivation}

Any single window's prediction, however confident, is a noisy
estimate: a brief muscle twitch, an eye movement unrelated to a
deliberate blink, or an artefact near a class boundary can produce
an isolated, spurious high-confidence prediction. Rather than
acting on every qualifying window individually, the system requires
a short run of \textbf{agreement} across consecutive predictions
before treating a class as a confirmed command.

\subsection{Design}

The \acrshort{fsm} maintains a simple run-length counter over
predictions that have already passed the activity gate and the
confidence threshold: it tracks the current candidate class and how
many consecutive qualifying windows have agreed with it. A command
is confirmed, and transmitted over \acrshort{uart}
(Section~\ref{sec:inf_uart}), only once \textbf{3 consecutive}
qualifying windows have produced the \emph{same} predicted class.
Figure~\ref{fig:fsm} illustrates the state machine.

\begin{figure}[htbp]
\centering
\begin{tikzpicture}[
    node distance=1.6cm,
    state/.style={circle, draw=black, fill=blue!10,
                  minimum size=1.6cm, align=center, font=\small},
    final/.style={circle, draw=black, fill=green!20, thick,
                  minimum size=1.6cm, align=center, font=\small},
    arrow/.style={-{Stealth[length=5pt]}, thick}
]
\node[state] (s0) {0 matches\\(idle)};
\node[state, right=2.4cm of s0] (s1) {1 match};
\node[state, right=2.4cm of s1] (s2) {2 matches};
\node[final, right=2.4cm of s2] (s3) {Confirmed\\emit UART};

\draw[arrow] (s0) to[bend left=25] node[above,
    font=\scriptsize]{qualifying pred.} (s1);
\draw[arrow] (s1) to[bend left=25] node[above,
    font=\scriptsize]{same class} (s2);
\draw[arrow] (s2) to[bend left=25] node[above,
    font=\scriptsize]{same class} (s3);
\draw[arrow] (s1) to[out=225,in=315] node[below,
    font=\scriptsize]{different class} (s0);
\draw[arrow] (s2) to[out=225,in=315] node[below,
    font=\scriptsize]{different class / restart at 1} (s0);
\draw[arrow] (s3) to[out=135,in=45,looseness=2.5]
    node[above, font=\scriptsize]{reset} (s0);
\end{tikzpicture}
\caption{The 3-consecutive-prediction confirmation FSM. Any
disagreement resets the run; a window that fails the activity gate
or confidence threshold (Sections~\ref{sec:inf_gate}--\ref{sec:inf_confidence})
does not advance or reset the counter, since it is treated as ``no
information'' rather than ``a competing class.''}
\label{fig:fsm}
\end{figure}

A window that fails the activity gate or the confidence threshold
is treated as a non-event rather than a competing prediction: it
neither advances the run-length counter nor resets it, so a brief
rest period between two windows supporting the same gesture does
not force the count back to zero. Only a \emph{qualifying}
prediction of a \emph{different} class resets the counter.

%% ----------------------------------------------------------
\section{Latency Budget}
\label{sec:inf_latency}
%% ----------------------------------------------------------

Measured end-to-end, per-window inference (feature extraction
through classification, for whichever branch a window is routed to)
takes \textbf{14\,ms}, dominated by the \acrshort{dtw} distance
computation in Stage~2a --- the most expensive single operation in
the pipeline, since it involves a Sakoe--Chiba-constrained dynamic
programming pass against every stored training exemplar
(Chapter~9). This is well under the 200\,ms window step
(Section~\ref{sec:inf_window}): at 14\,ms of compute per 200\,ms of
wall-clock step, the pipeline uses roughly 7\% of its available time
budget per window, leaving very substantial headroom before the
real-time constraint (finishing one window's classification before
the next window is due) becomes a concern. This margin is also what
makes the \acrshort{fsm}'s 3-consecutive-window confirmation
practical without perceptible added lag: at five predictions per
second, three consecutive agreeing predictions are available within
roughly 0.6\,s of a gesture's onset in the best case, well inside
what a user would perceive as a responsive control loop.

%% ----------------------------------------------------------
\section{UART Output (Contract A)}
\label{sec:inf_uart}
%% ----------------------------------------------------------

\subsection{Correcting an Outdated Description}

Earlier working notes describing this pipeline stage referred to
the classifier's output as going to ``Arduino serial'' --- a
description carried over from an earlier hardware iteration of the
project, before the current 3-node distributed architecture
(Chapter~4) was adopted. \textbf{This is out of date and is
corrected here explicitly:} the deployment target is not an
Arduino, and the link is not a generic serial print. The confirmed
command is transmitted as a single \code{lowercase\_snake\_case}
label, newline-terminated, over the Raspberry Pi~5's UART pins to
the ESP32's \code{Serial2} input at 115200 baud --- precisely
\textbf{Contract~A}, specified in full in Chapter~11. This chapter
and Chapter~11 describe the two ends of the same physical link: this
chapter produces the label; Chapter~11 describes how the ESP32
consumes it, looks it up in the user-editable mapping table, and
drives the servos.

\subsection{Label Emission}

On \acrshort{fsm} confirmation, the pipeline emits exactly the
class name the model predicted --- \code{single\_blink},
\code{double\_blink}, \code{clinch}, \code{bruxism\_left}, or
\code{bruxism\_right} --- as the UART line. As documented in
Chapter~11 (Section~11.4.2), the currently shipped ESP32 firmware
configuration does not yet recognise these exact label strings
(it ships with an older seven-label placeholder set); the
reconciled mapping table given there is what this pipeline's output
is designed against; that reconciliation is a firmware-side
configuration fix, not a change required in the inference pipeline
described in this chapter.

\subsection{Command Line Interface}

The production inference script, \code{inference.py}, accepts a
\code{--model} flag selecting which trained architecture to run
live: \code{dtw} (the final production model described throughout
this chapter), \code{hybrid} (Riemannian muscle branch only, hand-crafted
blink features, Chapter~9), \code{hierarchical}, or \code{flat}. This
was primarily a development and A/B-testing convenience during the
architecture progression documented in Chapter~9, allowing each
intermediate architecture to be exercised live on the actual hand
before its offline evaluation numbers were finalised; \code{dtw} is
the flag used for all deployed operation.

%% ----------------------------------------------------------
\section{Offline vs.\ Real-Time: A Consolidated Comparison}
\label{sec:inf_compare}
%% ----------------------------------------------------------

Table~\ref{tab:offline_vs_realtime} consolidates every point in this
report where the offline (training/evaluation) pipeline and the
real-time (deployment) pipeline intentionally differ, since the two
are easy to conflate and have been addressed piecemeal across
several chapters.

\begin{table}[htbp]
\centering
\caption{Every intentional difference between the offline and
real-time pipelines. Everything not listed here (feature
extraction, the trained model itself, channel configuration) is
identical between the two.}
\label{tab:offline_vs_realtime}
\renewcommand{\arraystretch}{1.3}
\begin{tabular}{p{2.8cm} p{4.9cm} p{4.9cm}}
\toprule
\textbf{Aspect} & \textbf{Offline (training/eval)} &
\textbf{Real-time (deployment)} \\
\midrule
Window step & 0.5\,s (75\% overlap; Chapter~7) & 0.2\,s (5
    predictions/s; Section~\ref{sec:inf_window}) \\
Activity gate & Recording-relative,
    \code{gate\_mask(frac=0.5)}, 90th-percentile-based
    (Chapter~7) & Fixed 20\,\si{\micro\volt} \acrshort{rms}
    threshold (Section~\ref{sec:inf_gate}) \\
Evaluation unit & Independent epoch, scored against ground truth
    (Chapter~9) & Continuous stream, no ground truth, confirmed
    via \acrshort{fsm} (Section~\ref{sec:inf_fsm}) \\
Output & Predicted label + true label, for metric computation
    (Chapter~10) & Confirmed label only, transmitted over
    \acrshort{uart} (Section~\ref{sec:inf_uart}) \\
\bottomrule
\end{tabular}
\end{table}

%% ----------------------------------------------------------
\section{Chapter Summary}
\label{sec:inf_summary}
%% ----------------------------------------------------------

This chapter described the deployment of the offline-trained
\acrshort{dtw} hybrid model (Chapter~9) as a real-time classifier.
A rolling 2.0\,s window, re-evaluated every 0.2\,s, is passed
through a \textbf{fixed} 20\,\si{\micro\volt} activity gate ---
distinct from, and for good reason not identical to, the
recording-relative gate used to build the training dataset
(Chapter~7) --- before being classified by the unmodified Chapter~9
model and gated a second time by a 0.60 confidence threshold. A
3-consecutive-prediction Finite State Machine suppresses isolated,
noisy predictions before any command is confirmed. Measured
per-window latency (14\,ms, dominated by the \acrshort{dtw}
distance computation) leaves substantial headroom under the 200\,ms
step, so the pipeline comfortably meets real-time throughput
requirements. Confirmed commands are transmitted as
\code{lowercase\_snake\_case} class labels over \acrshort{uart} to
the ESP32 via Contract~A (Chapter~11) --- explicitly correcting an
outdated ``Arduino serial'' description carried over from an
earlier hardware iteration of the project (Chapter~4). A
consolidated table cross-referenced every intentional offline/real-time
difference discussed across this and earlier chapters. The
following chapter discusses the project's results, limitations, and
broader implications.